\subsection*{6.4 Divergence Theorem}

\begin{conceptbox}
    \begin{equation*}
        \iiint_V (\nabla \cdot \mathbf{F})\, dV = \iint_S \mathbf{F}\cdot d\mathbf{S}
    \end{equation*}
\end{conceptbox}

\begin{itemize}
    \item Volume integral $\rightarrow$ total source (how much the field is “generated” inside)
    \item Surface integral $\rightarrow$ total flux (how much exits the boundary)
\end{itemize}

The Divergence Theorem states that the total outward flow (flux) of a vector field across a closed surface is equal to the total amount of “source” inside the volume.

Instead of computing a difficult surface integral directly, we can:
\begin{itemize}
    \item compute the divergence (a scalar function),
    \item integrate it over the volume,
    \item and obtain the same result.
\end{itemize}

\[
    \boxed{\text{Flux out of surface = Total expansion inside volume}}
\]

\begin{examplebox}
    Compute the flux of $\mathbf{F}=\langle x^3,y^3,z^3\rangle$ over the sphere $x^2+y^2+z^2\le1$.

    \textbf{Step 1: Understand the problem}

    We are asked to compute:
    \[
        \iint_S \mathbf{F}\cdot d\mathbf{S}
    \]

    where $S$ is the surface of the unit sphere.

    Direct computation would require parametrizing the sphere and computing a surface integral, which is complicated.

    Instead, we apply the Divergence Theorem:
    \[
        \iint_S \mathbf{F}\cdot d\mathbf{S}
        =
        \iiint_V (\nabla \cdot \mathbf{F})\, dV
    \]

    \textbf{Step 2: Compute the divergence}

    Given:
    \[
        \mathbf{F} = \langle x^3,y^3,z^3\rangle
    \]

    Compute divergence:
    \begin{align*}
        \nabla \cdot \mathbf{F}
         & = \frac{\partial}{\partial x}(x^3) + \frac{\partial}{\partial y}(y^3) + \frac{\partial}{\partial z}(z^3) \\
         & = 3x^2 + 3y^2 + 3z^2
    \end{align*}

    Factor:
    \begin{align*}
        \nabla \cdot \mathbf{F} = 3(x^2+y^2+z^2)
    \end{align*}

    Since:
    \[
        x^2+y^2+z^2 = r^2
    \]

    we get:
    \[
        \nabla \cdot \mathbf{F} = 3r^2
    \]

    \textbf{Step 3: Convert to spherical coordinates}

    The region is a sphere:
    \[
        x^2+y^2+z^2 \le 1
    \]

    So we use spherical coordinates:

    \begin{align*}
        x & = r\sin\theta\cos\phi \\
        y & = r\sin\theta\sin\phi \\
        z & = r\cos\theta
    \end{align*}

    Volume element:
    \[
        dV = r^2\sin\theta\,dr\,d\theta\,d\phi
    \]

    \textbf{Step 4: Substitute into the integral}

    \begin{align*}
        \iiint_V (\nabla \cdot \mathbf{F})\, dV
         & = \iiint_V 3r^2\, dV                                                           \\
         & = \int_0^{2\pi}\int_0^\pi\int_0^1 3r^2 \cdot r^2\sin\theta\,dr\,d\theta\,d\phi \\
         & = \int_0^{2\pi}\int_0^\pi\int_0^1 3r^4\sin\theta\,dr\,d\theta\,d\phi
    \end{align*}
    \textbf{Step 5: Separate the integrals}

    Because the variables are independent:
    \begin{align*}
        = 3 \left(\int_0^1 r^4\,dr\right)
        \left(\int_0^\pi \sin\theta\,d\theta\right)
        \left(\int_0^{2\pi} d\phi\right)
    \end{align*}

    \textbf{Step 6: Evaluate each integral}

    \textbf{Radial part:}
    \begin{align*}
        \int_0^1 r^4\,dr = \left[\frac{r^5}{5}\right]_0^1 = \frac{1}{5}
    \end{align*}

    \textbf{Angular ($\theta$):}
    \begin{align*}
        \int_0^\pi \sin\theta\,d\theta = [-\cos\theta]_0^\pi = 2
    \end{align*}

    \textbf{Azimuthal ($\phi$):}
    \begin{align*}
        \int_0^{2\pi} d\phi = 2\pi
    \end{align*}

    \textbf{Step 7: Multiply results}

    \begin{align*}
        3 \cdot \frac{1}{5} \cdot 2 \cdot 2\pi = \frac{12\pi}{5}
    \end{align*}

    \textbf{Final Answer:}
    \[
        \boxed{\frac{12\pi}{5}}
    \]

\end{examplebox}